\section{Cryptographic Properties}
\label{sec:cryptographic-properties}

\subsection{Threat model}

The verifier operates in the following threat model:

\begin{itemize}[leftmargin=2em,nosep]
\item \textbf{Adversary capability.} Probabilistic polynomial-time
  Turing machine. PQ adversary is a polynomial-time quantum Turing
  machine; the verifier's PQ claim is against this stronger
  adversary class.
\item \textbf{Adversary access.} The adversary chooses calldata
  (proof bytes and public inputs) and submits it to a node's EVM.
  The adversary observes the precompile's return tuple but not the
  validator's internal state.
\item \textbf{Honest verifier.} The validator runs the unmodified
  precompile + backend; we assume no fault injection on the
  validator side.
\item \textbf{Random oracle for cSHAKE256.} cSHAKE256 is modeled as
  a random oracle in the security reduction. This is the standard
  assumption for FRI's Fiat-Shamir compilation; the modeled
  primitive is the cSHAKE256 customizable XOF as standardized in
  FIPS-202 / SP-800-185.
\end{itemize}

There are \emph{no secret keys} on the verifier side. STARK proofs
authenticate public statements; the verifier's secrets-side surface
is empty. The CT discussion in Section~\ref{sec:constant-time}
therefore concerns soundness leakage (which validation stage
rejected) rather than confidentiality leakage (no confidentiality
to leak).

\subsection{Soundness}

\subsubsection{Knowledge soundness}

Under the random-oracle model for cSHAKE256, FRI is a knowledge-sound
interactive oracle proof~\cite{frieprint,starkeprint}. Knowledge
soundness extracts a witness from any prover that convinces the
verifier with probability above the soundness error.

Concretely, for the strict-PQ Plonky3-fork configuration used by the
audited backend crate, the knowledge-soundness statement is:

\begin{quote}
\emph{There exists a polynomial-time extractor $E^{P}$ such that for
every PPT prover $P$ that wins the FRI knowledge-soundness game with
probability $\epsilon$, the extractor outputs a witness for the
underlying R1CS / Plonkish instance with probability
$\ge \epsilon - 2^{-100}$.}
\end{quote}

The $2^{-100}$ figure is the \emph{soundness error}, derived from
the FRI parameter set: query count $q$, blowup factor $\rho$,
folding factor $f$, and field size $|\mathbb{F}|$. The Plonky3-fork
parameter set is documented in
\texttt{lux/p3q/crates/p3q-verifier/src/params.rs}; LP-221 pins
the verifier-conformance, not the prover parameter schedule (the
prover may choose stronger parameters but cannot choose weaker
ones, because the verifier hard-codes its expected schedule).

\subsubsection{Soundness error}

Per the parameter schedule used by the audited backend crate, the
soundness error is $\le 2^{-100}$ in the random-oracle model. This
is the conventional STARK target~\cite{starkeprint}: it bounds the
probability that a cheating prover convinces the verifier of a
false statement.

\subsubsection{Completeness}

Honest prover, honest verifier: the verifier accepts with
probability~$1$. There is no statistical-zero-knowledge slack at
the verify side; the prover may have ZK slack (depending on whether
the Plonky3-fork hiding mode is enabled), but the verifier's
acceptance probability on an honestly-generated proof is exactly~$1$.

\subsubsection{Public-coin / Fiat-Shamir}

FRI is public-coin~\cite{frieprint}; the Fiat-Shamir transform
with cSHAKE256 as the random oracle compiles it to a
non-interactive argument. The transcript hashing
domain-separates rounds; cSHAKE256's customization string carries
the per-round domain tag.

\subsection{Post-quantum assumption}

The single cryptographic assumption is \textbf{collision resistance
of cSHAKE256}, modeled as a random oracle in the Fiat-Shamir
reduction. There is no discrete-log, no pairing, no factoring, no
LWE, no SVP, no SIDH-style isogeny assumption.

cSHAKE256~\cite{fips202,sp800185} inherits the SHA-3 / Keccak-$f[1600]$
design. The conjectured security strength is:

\begin{table}[h]
\centering
\small
\begin{tabular}{@{}llr@{}}
\toprule
\textbf{Adversary} & \textbf{Best attack} & \textbf{Conjectured cost} \\
\midrule
Classical & generic preimage / collision (birthday) & $\ge 2^{256}$ preimage; $\ge 2^{128}$ collision \\
Quantum   & Grover preimage; BHT collision           & $\ge 2^{128}$ preimage; $\ge 2^{85}$ collision \\
\bottomrule
\end{tabular}
\caption{cSHAKE256 security strength against classical and quantum
adversaries at the 256-bit output length. Adequate for
$\le 2^{-100}$ soundness error at the FRI parameter choices the
backend uses; the soundness reduction is loose by $\sim$$30$~bits
of margin against the Grover-bounded quantum attacker, which is
the standard PQ-STARK safety budget.}
\label{tab:cshake-strength}
\end{table}

\subsubsection{Why not BLAKE3 or Poseidon}

LP-221 picks cSHAKE256 over alternatives for three reasons:

\begin{enumerate}[leftmargin=2em,nosep]
\item \textbf{Standardization.} cSHAKE256 is FIPS-202~/
  SP-800-185~\cite{fips202,sp800185}. NIST standardization gates
  conformance test vectors, third-party crypto-library
  implementations, and audit checklists. BLAKE3 is widely deployed
  but not FIPS-standardized; Poseidon is SNARK-friendly but its
  collision-resistance under standard-model assumptions is the
  subject of ongoing analysis.
\item \textbf{Customization string.} cSHAKE256's customization
  string ($S$ parameter) is purpose-built for domain separation
  between protocol contexts. The Plonky3-fork uses per-round
  customization strings to bind Fiat-Shamir transcripts;
  reusing the same primitive with different $S$ values is
  cleaner than concatenating tag-prefixes onto a non-customizable
  hash function.
\item \textbf{Audit surface uniformity.} The Quasar stack uses
  Keccak-$f[1600]$ as the hash primitive at multiple layers
  (cSHAKE256 here, SLH-DSA inside Magnetar~\cite{fips205}, EVM
  KECCAK256 at the precompile boundary). Co-locating the hash
  primitive reduces the audit surface and amortizes the
  side-channel hardening across multiple protocol layers.
\end{enumerate}

\subsection{Domain separation}

The on-chain wire format carries a fixed leading
\texttt{MagicHeader = "P3Q1"} that ties the precompile to the
strict-PQ Plonky3 fork. Within the proof bytes, the Plonky3-fork
verifier applies its own domain separation via cSHAKE256
customization strings on each Fiat-Shamir round; that domain
separation is internal to the backend crate and is out of scope of
LP-221's wire-format contract.

The two-layer domain separation (wire-bytes magic header at the
EVM boundary $+$ Fiat-Shamir customization inside the backend) is
deliberate. The EVM boundary's magic header is a coarse
type-discriminator: it answers ``is this a STARK-FRI strict-PQ v1
proof or something else.'' The backend's customization strings are
fine-grained: they answer ``which round of which proof is this
hash invocation for.'' Conflating the two is an anti-pattern
(``omit the magic header because the backend will catch it later'')
because it shifts a cheap, structural check into an expensive,
cryptographic one.

\subsection{What is NOT claimed}

\begin{itemize}[leftmargin=2em,nosep]
\item \textbf{Zero-knowledge.} The Plonky3-fork verifier as deployed
  does NOT carry a zero-knowledge property at the LP-221 boundary.
  Public inputs (\texttt{prev\_root}, \texttt{new\_root},
  \texttt{batch\_hash}) are explicitly public. A future LP can
  graft a ZK-extension by activating Plonky3's hiding commitment
  mode and bumping the version byte. Today, LP-221 verifies
  computational integrity, not zero-knowledge proof-of-knowledge.
\item \textbf{Universal setup.} Not applicable; FRI is transparent
  (no trusted setup of any kind). The Plonky3 fork inherits this
  property; LP-221 does not introduce any setup-requiring
  component.
\item \textbf{Recursion.} The Plonky3-fork verifier as exposed at
  LP-221 does not recurse over itself on-chain; recursion happens
  off-chain in the prover and surfaces to the EVM as a single
  STARK proof. A future recursive variant would land at a separate
  version byte (and likely separate magic header) under the same
  slot.
\item \textbf{Quantitative GPU verify performance.} LP-221 does NOT
  claim measured GPU throughput. The CPU verify cost of
  \SIrange{5}{10}{\milli\second} on a modern x86\_64 core is the
  measured number; the projected $\sim$\SI{50}{\micro\second} per
  verify on Blackwell sm\_120 is a projection from LP-203 kernel
  micro-benchmarks~\cite{lp203}, not a measurement. The gas
  schedule is sized to the CPU-class cost.
\end{itemize}
